% =============================================================================
% Minimal siunitx replacement.
%
% Loaded by main.tex ONLY when siunitx.sty is unavailable (on Debian/Ubuntu
% siunitx ships in texlive-science, not texlive-latex-extra). It implements
% exactly the macros this paper uses: \SI, \si, \num, \SIrange, \numrange and
% the unit names \meter \second \gram \newton \radian \hertz \percent \degree
% \g \kilo \milli \per \squared.
%
% All macros are declared robust: table captions are moving arguments
% (written to the .lot file), and an expandable conditional breaks there.
% Output matches siunitx closely; the one visible difference is that
% \SIrange/\numrange print an en-dash here where siunitx prints " to ".
% Install texlive-science to get the real package.
% =============================================================================

\makeatletter

% Separator state. Armed (true) once a base unit has been emitted; disarmed
% after a prefix and after \per, so "kilo newton per radian" sets no space
% between k and N, nor around the solidus.
\newif\ifsiu@prev
\DeclareRobustCommand{\siu@sep}{\ifsiu@prev\,\fi}
% \text{} (amsmath) so a unit typesets upright in BOTH text and math mode --
% these macros are routinely used inside $...$, where a bare `m` would come out
% as an italic math variable and `--` as two minus signs.
\DeclareRobustCommand{\siu@unit}[1]{\siu@sep\text{#1}\siu@prevtrue}  % base unit
\DeclareRobustCommand{\siu@pre}[1]{\siu@sep\text{#1}\siu@prevfalse}  % SI prefix
\DeclareRobustCommand{\siu@div}[1]{\text{#1}\siu@prevfalse}          % solidus

% Default (outside \si) meanings, so the unit names are always defined and
% \renewcommand below always has a target.
\DeclareRobustCommand{\meter}{m}
\DeclareRobustCommand{\second}{s}
\DeclareRobustCommand{\gram}{g}
\DeclareRobustCommand{\newton}{N}
\DeclareRobustCommand{\radian}{rad}
\DeclareRobustCommand{\hertz}{Hz}
\DeclareRobustCommand{\percent}{\%}
\DeclareRobustCommand{\g}{\textit{g}}
\DeclareRobustCommand{\kilo}{k}
\DeclareRobustCommand{\milli}{m}
\DeclareRobustCommand{\per}{/}
\DeclareRobustCommand{\squared}{\ensuremath{^{2}}}
\providecommand{\degree}{\ensuremath{^{\circ}}}

\DeclareRobustCommand{\si}[1]{\begingroup
  \siu@prevfalse
  \renewcommand{\meter}{\siu@unit{m}}%
  \renewcommand{\second}{\siu@unit{s}}%
  \renewcommand{\gram}{\siu@unit{g}}%
  \renewcommand{\newton}{\siu@unit{N}}%
  \renewcommand{\radian}{\siu@unit{rad}}%
  \renewcommand{\hertz}{\siu@unit{Hz}}%
  \renewcommand{\percent}{\siu@unit{\%}}%
  \renewcommand{\degree}{\siu@unit{\ensuremath{^{\circ}}}}%
  \renewcommand{\g}{\siu@unit{\textit{g}}}%
  \renewcommand{\kilo}{\siu@pre{k}}%
  \renewcommand{\milli}{\siu@pre{m}}%
  \renewcommand{\per}{\siu@div{/}}%
  \renewcommand{\squared}{\ensuremath{^{2}}}%
  #1\endgroup}

\DeclareRobustCommand{\num}[1]{#1}
\DeclareRobustCommand{\numrange}[2]{#1\text{--}#2}
\DeclareRobustCommand{\SI}[2]{#1\,\si{#2}}
\DeclareRobustCommand{\SIrange}[3]{#1\text{--}#2\,\si{#3}}

\makeatother
